\documentclass{report}
\usepackage{apalike}
\usepackage{comment}
\usepackage[utf8]{inputenc}
\usepackage[spanish]{babel}
\usepackage{amsthm}
\usepackage{hyperref}
\usepackage{booktabs}
\usepackage{amsmath}
\usepackage{amsfonts}
\usepackage{amssymb}
\usepackage{graphicx}
\usepackage{float}
\usepackage{url}
\usepackage[left=2cm,right=2cm,top=2cm,bottom=2cm]{geometry}
\usepackage{fancyhdr}
\usepackage{array}
\usepackage{longtable}
\usepackage{enumitem}

\newtheorem{definition}{Definition}[section]
\newtheorem{example}{Example}[section]
\newtheorem{theorem}{Theorem}[section]

\pagestyle{fancy}
\fancyhf{}
\fancyhead[L]{\includegraphics[width=2cm]{Logo_Loyola.png}}
\fancyhead[C]{\textbf{Desarrollo Alumni}}
\fancyhead[R]{\today}
\fancyfoot[C]{\thepage}

\graphicspath{{Images/}}

\title{\textbf{Práctica: Desarrollo Alumni}}
\author{
    Fernando Cáceres Plaza\\
    Francisco Manuel Moya Sánchez\\
    Diana Luque Pérez\\
    \textbf{PDM 2025-2026} \\
    \textbf{Institución: Universidad Loyola Andalucía}
}
\date{}

\renewcommand{\thesection}{\arabic{section}}
\renewcommand{\thesubsection}{\arabic{section}.\arabic{subsection}}

\newenvironment{usecasetable}{%
    \par\noindent
    \begin{longtable}{|p{0.2\textwidth}|p{0.7\textwidth}|}
    \hline
}{%
    \end{longtable}
    \newpage
}

\newenvironment{flowtable}[1]{%
    \begin{tabular}{|l|p{#1}|}
    \hline
    \textbf{Paso} & \textbf{Acción} \\
    \hline
}{%
    \end{tabular}
}

\newcommand{\dependencies}[1]{%
    \begin{itemize}[leftmargin=*, nosep]
        #1
    \end{itemize}
}

\begin{document}
\maketitle
\newpage
\pagenumbering{arabic}
\setcounter{page}{1}
\tableofcontents
\newpage

\section{Problema real y problema técnico}
\subsection{Problema real}
Somos una empresa de desarrollo de software y nos ha contactado la universidad Loyola para desarrollar un sistema que le permita poner en contacto a antiguos alumnos de la universidad, también quiere que a partir de este sistema podemos organizar eventos entre empresas y estos antiguos alumnos. Debemos empezar desde cero y teniendo como referencia algunas páginas web de alumnis de otras universidades crear una aplicacion que sea sostenible en el tiempo y con posibilidad de crecimiento. Una vez finalizado también deberemos llevar su mantenimiento e ir actualizando el programa.

\subsection{Problema técnico (PDS)}
Para definir el problema desde el punto de vista de los desarrolladores, se utilizará la metodología PDS (Product Design Specification).

\subsubsection{Funcionamiento}
\begin{itemize}
    \item{\underline{Gestión de Usuarios, Roles y Accesos:}} El sistema permitirá una gestión integral. Roles:
        \begin{itemize}
            \item{Alumni:} Modifican datos, buscan otros alumni, proponen eventos y se inscriben.
            \item{Administradores:} Control total, validación de propuestas y auditoría.
            \item{PTGAS:} Gestionan y evalúan eventos y propuestas.
            \item{PDI:} Buscan alumni y proponen eventos.
        \end{itemize}
    \item{\underline{Filtros de búsqueda:}} Consulta avanzada por nombre, titulación, campus, trabajo, hobbies.
    \item{\underline{Gestión de eventos:}} Flujo de propuestas por parte de Alumni/PDI, evaluadas por PTGAS/Admin, y posterior publicación.
\end{itemize}

\subsubsection{Entorno}
\begin{itemize}
    \item{\underline{Entorno software:}} Compatible con Salesforce (integración futura). La solución se entrega como \textbf{App Android nativa} (Java 11, Retrofit, OkHttp) como cliente principal contra un backend Java EE sobre \textbf{Apache Tomcat 9}; la opción PWA queda como entorno alternativo / web a futuro.
    \item{\underline{Entorno hardware:}} Libre elección. Salesforce gestionará relaciones CRM en el futuro.
    \item{\underline{Entorno de usuario:}} Interfaz móvil intuitiva sobre dispositivos Android.
\end{itemize}

\subsubsection{Vida esperada y Ciclo de mantenimiento}
Operativo a la mayor brevedad. Mantenimiento continuo con revisiones semanales y soporte a la app Android nativa y, en el futuro, a la PWA.

\subsubsection{Competencia}
Ejemplos: Alumni Global IE, ESADE Alumni, Universidad de Navarra.

\subsubsection{Aspecto externo}
\begin{itemize}
    \item{\underline{Soporte y acceso:}} La solución entregada es una \textbf{aplicación Android nativa} instalable como APK (compileSdk 36, minSdk 24). Hará uso de sensores y servicios del dispositivo mediante el SDK de Android (intents implícitos para SMS, etc.). Una futura PWA podrá replicar la experiencia en navegadores modernos.
    \item{\underline{Documentación e Interfaz:}} Manual técnico, diseño corporativo Loyola, accesibilidad básica.
\end{itemize}

\subsubsection{Estandarización, Calidad y Tareas}
Metodología SCRUM con Sprints. Pruebas funcionales, de seguridad y rendimiento.

\subsubsection{Seguridad}
Acceso estricto por roles, privacidad por defecto (RGPD) y minimización de datos.

\section{Especificación de Requisitos}
\subsection{Requisitos Generales}
\begin{itemize}
    \item \textbf{RG-1.} Gestión de Usuarios y Accesos
    \item \textbf{RG-2.} Búsqueda y Networking de Alumni
    \item \textbf{RG-3.} Gestión de Eventos y Actividades
    \item \textbf{RG-4.} Seguridad y Privacidad del Dato
    \item \textbf{RG-5.} Integración con Ecosistema (Salesforce) [\textit{conceptual, futura}]
    \item \textbf{RG-6.} Experiencia Móvil (\textbf{Android nativa}) y Usabilidad
    \item \textbf{RG-7.} Mantenimiento, Auditoría y Escalabilidad
\end{itemize}

\subsection{Requisitos Específicos}
\subsubsection{Requisitos Funcionales}
\begin{itemize}
    \item \textbf{RF-1.} Los usuarios accederán mediante credenciales tras validación de rol.
    \item \textbf{RF-2.} El sistema permitirá la activación de cuenta mediante un enlace temporal seguro enviado por correo, exigiendo el cambio de contraseña en el primer acceso y aplicando política de caducidad.
    \item \textbf{RF-3.} Los roles autorizados (Alumni, PDI, PTGAS, Admin) podrán usar búsquedas simples y avanzadas aplicando filtros cruzados.
    \item \textbf{RF-4.} Los Alumni y PDI podrán registrar propuestas de eventos y actividades que quedarán en estado pendiente.
    \item \textbf{RF-5.} El PTGAS podrá modificar su propio perfil y buscar perfiles de Alumni.
    \item \textbf{RF-6.} Los administradores podrán crear, modificar, suspender (baja lógica) y cambiar roles de cuentas.
    \item \textbf{RF-7.} El PTGAS y Administradores podrán revisar, evaluar (aprobar o rechazar) las propuestas de eventos y actividades.
    \item \textbf{RF-8.} Los administradores y PTGAS podrán publicar eventos y actividades aprobados, modificarlos o cancelarlos.
    \item \textbf{RF-9.} Los administradores tendrán acceso a un panel de control con métricas de usuarios, eventos, propuestas pendientes, inscripciones y auditoría de acciones.
    \item \textbf{RF-10.} Los Alumni y PDI podrán inscribirse o cancelar su inscripción a eventos/actividades aprobadas, siempre que existan plazas y el plazo esté abierto.
    \item \textbf{RF-11.} El alumni podrá editar su propio perfil, gestionando explícitamente sus preferencias de visibilidad por campo.
    \item \textbf{RF-12.} El sistema ejecutará un proceso automático para notificar a cuentas con inactividad superior a un año.
    \item \textbf{RF-13.} La visibilidad de los perfiles estará estrictamente condicionada a la autenticación, la autorización del rol que busca y las preferencias de privacidad del alumni objetivo.
    \item \textbf{RF-14.} Tras crear una cuenta de usuario, el sistema permitirá al administrador lanzar un aviso de bienvenida al nuevo usuario por SMS mediante un \textbf{intent implícito} del sistema operativo Android (\texttt{Intent.ACTION\_SENDTO} con URI \texttt{smsto:}), que delegará la entrega en la app de mensajería predeterminada del dispositivo. El mensaje no incluirá credenciales. Dependencias: RG-1, RF-6.
\end{itemize}

\subsubsection{Requisitos de Información}
\begin{itemize}
    \item \textbf{RI-1.} Perfil Alumni: nombre, titulación, promoción, campus, teléfono, email.
    \item \textbf{RI-2.} Trabajo y Hobbies: descripciones, empresa, ciudad, experiencia.
    \item \textbf{RI-3.} Credenciales y Roles: usuario, hash de contraseña segura, rol.
    \item \textbf{RI-4.} Eventos y Actividades: nombre, tipo de recurso (evento o actividad), fechas, aforo, ponentes/responsables, estado y ubicación.
    \item \textbf{RI-5.} Inscripciones: registro de alumni/pdi, recurso inscrito (evento o actividad), fecha de inscripción, asistencia y estado de cancelación.
    \item \textbf{RI-6.} Propuestas de Eventos y Actividades: tipo de propuesta, estado (pendiente, aprobada, rechazada), motivo, solicitante, fecha de envío y decisión asociada.
    \item \textbf{RI-7.} Perfil PDI/PTGAS: nombre, correo, área o departamento.
    \item \textbf{RI-8.} Auditoría e Historial de Cambios: fecha, actor, acción sensible realizada, entidad afectada, resultado y valores relevantes antes/después cuando proceda.
    \item \textbf{RI-9.} Privacidad: matriz de preferencias de visibilidad campo a campo por Alumni.
    \item \textbf{RI-10.} Integración con Salesforce: tokens de acceso, fecha de última sincronización, mapeo de IDs externos y estado de conexión. [\textit{conceptual, integración futura}]
    \item \textbf{RI-11.} Notificaciones y Recordatorios: tipo de notificación (\texttt{INACTIVIDAD}, \texttt{PROPUESTA\_RESUELTA}, \texttt{EVENTO\_RECORDATORIO}, \texttt{SISTEMA}), destinatario, fecha de envío y estado (leído/no leído). \textbf{Implementado} mediante la tabla \texttt{notificacion}, la clase DAO \texttt{ManagerNotificaciones} y el endpoint \texttt{/NotificacionServlet}.
\end{itemize}

\subsubsection{Reglas de Negocio}
\begin{itemize}
    \item \textbf{RN-1.} Solo los usuarios autenticados y pertenecientes a colectivos autorizados podrán acceder a la plataforma.
    \item \textbf{RN-2.} Cada alumni controlará la visibilidad de los datos personales de su perfil.
    \item \textbf{RN-3.} Por defecto, los datos personales no serán visibles salvo los mínimos necesarios para identificar el perfil.
    \item \textbf{RN-4.} Las búsquedas solo mostrarán la información permitida por las preferencias de privacidad del alumni.
    \item \textbf{RN-5.} Las propuestas de eventos o actividades de Alumni/PDI deberán quedar en estado pendiente hasta su revisión.
    \item \textbf{RN-6.} Solo los eventos o actividades aprobados serán visibles e inscribibles.
    \item \textbf{RN-7.} No se permitirá la inscripción fuera del plazo definido.
    \item \textbf{RN-8.} No se permitirá superar el número de plazas disponibles.
    \item \textbf{RN-9.} Un usuario no podrá inscribirse dos veces al mismo evento o actividad.
    \item \textbf{RN-10.} La eliminación de una cuenta deberá ser una baja lógica o anonimizar los datos históricos para conservar la trazabilidad (RGPD).
    \item \textbf{RN-11.} Toda acción administrativa sensible deberá quedar registrada en auditoría.
    \item \textbf{RN-12.} Las cuentas inactivas durante un año recibirán un recordatorio según la política definida. El proceso lo ejecuta automáticamente el \textit{scheduler} interno \texttt{JobsContextListener} cada 24 horas (ventana 03:00 local), además del disparo manual del endpoint correspondiente.
\end{itemize}

\subsubsection{Requisitos No Funcionales}
\begin{itemize}
    \item \textbf{RNF-1:} La interfaz móvil se entrega como \textbf{app Android nativa} (compileSdk 36, minSdk 24, Material 3). Una futura PWA basada en HTML5, CSS3 y Service Workers complementará el acceso desde navegadores, garantizando una puntuación mínima de usabilidad y rendimiento del \textbf{90\%} en Google Lighthouse cuando se entregue.
    \item \textbf{RNF-2:} El tiempo máximo de respuesta del backend para cualquier petición GET de búsqueda o carga de perfil no superará los \textbf{2.0 segundos} bajo una carga normal de operación de hasta 100 usuarios concurrentes simultáneos.
    \item \textbf{RNF-3:} La plataforma garantizará una disponibilidad operativa mínima del \textbf{99.5\%} medida anualmente, excluyendo de este cómputo las ventanas de mantenimiento técnico programadas que se realizarán semanalmente en horario valle (domingos de 02:00 a 05:00 AM).
    \item \textbf{RNF-4:} El sistema ejecutará un proceso automatizado de copia de seguridad incremental diariamente a las 03:00 AM, y una copia de seguridad completa de forma semanal cada domingo, almacenada de forma redundante y encriptada en un servidor secundario aislado. \textbf{Implementado} mediante los scripts \texttt{backup\_completo.bat}, \texttt{backup\_incremental.bat}, \texttt{programar\_tareas.bat} y \texttt{restaurar\_copia.bat} ubicados en la carpeta \texttt{alumni/backup/}, basados en \texttt{mysqldump}, \texttt{openssl} (AES-256) y \texttt{schtasks} de Windows; las copias se duplican en dos ubicaciones (primaria y secundaria) con retención configurable.
    \item \textbf{RNF-5:} Todos los controladores (Servlets) interceptarán las excepciones globales de la aplicación utilizando la capa \texttt{ServletJsonUtil}. El backend responderá obligatoriamente bajo una estructura JSON estandarizada con la clave \texttt{\{"success": false, error: "mensaje"\}} acompañada del código de estado HTTP adecuado (400 para datos inválidos, 401 para fallos de sesión, 403 por falta de rol, 404 para recursos inexistentes y 500 para errores internos).
    \item \textbf{RNF-6:} El diseño visual del frontend cumplirá estrictamente con las pautas de accesibilidad para el contenido web \textbf{WCAG 2.1 en su nivel de conformidad AA}, garantizando contrastes mínimos de texto de 4.5:1 y la navegabilidad completa mediante lectura de pantalla y teclado.
    \item \textbf{RNF-7:} Toda la comunicación entre el cliente (app Android, futura PWA) y el servidor Tomcat backend viajará obligatoriamente cifrada utilizando el protocolo criptográfico \textbf{TLS 1.3 (HTTPS)}, rechazando conexiones no seguras en el puerto 80.
    \item \textbf{RNF-8:} La integración con Salesforce deberá realizarse mediante un servicio de sincronización documentado, definiendo datos intercambiados, autenticación, periodicidad, tratamiento de errores y mapeo de identificadores externos. [\textit{conceptual, integración futura}]
    \item \textbf{RNF-9:} El backend deberá estar organizado por capas diferenciadas (controladores, servicios/managers, acceso a datos y modelo de dominio), utilizando persistencia relacional mediante MariaDB/MySQL y permitiendo evolucionar hacia una arquitectura más robusta sin modificar la API pública.
    \item \textbf{RNF-10:} Todos los endpoints deberán comprobar autenticación, rol y permisos antes de ejecutar cualquier operación, garantizando privacidad por defecto y minimización de datos personales en las respuestas.
\end{itemize}

\subsection{Matrices de Trazabilidad}
\subsubsection*{Matriz de Trazabilidad de Implementación Completa (RF a Backend)}
\begin{center}
\begin{longtable}{|p{0.12\textwidth}|p{0.15\textwidth}|p{0.30\textwidth}|p{0.25\textwidth}|p{0.10\textwidth}|}
\hline
\textbf{RF} & \textbf{Caso Uso} & \textbf{Servicio Backend (Endpoint)} & \textbf{Clase / Controlador} & \textbf{Estado} \\ \hline
\endfirsthead
\hline
\textbf{RF} & \textbf{Caso Uso} & \textbf{Servicio Backend (Endpoint)} & \textbf{Clase / Controlador} & \textbf{Estado} \\ \hline
\endhead
RF-1 Login & CU-001 & \texttt{POST /LoginServlet} & \texttt{LoginServlet / ManagerUsuarios} & Impl. \\ \hline
RF-2 Activar & CU-009 & \texttt{POST /ActivarCuentaServlet} & \texttt{ActivarCuentaServlet / ManagerUsuarios} & Parcial \\ \hline
RF-3 Buscar & CU-002, CU-008 & \texttt{GET /BuscarAlumniServlet} & \texttt{BuscarAlumniServlet / AlumniManager} & Impl. \\ \hline
RF-4 Proponer & CU-007 & \texttt{POST /PropuestaServlet} & \texttt{PropuestaServlet / ManagerPropuestas} & Impl. \\ \hline
RF-5 PTGAS & CU-003, CU-008 & \texttt{POST /PerfilServlet} \& \texttt{GET /BuscarAlumniServlet} & \texttt{PerfilServlet / AlumniManager} & Impl. \\ \hline
RF-6 Admin & CU-006 & \texttt{POST/PUT /UsuarioAdminServlet} & \texttt{UsuarioAdminServlet / ManagerUsuarios} & Impl. \\ \hline
RF-7 Evaluar & CU-010 & \texttt{PUT /PropuestaServlet} & \texttt{PropuestaServlet / ManagerPropuestas} & Impl. \\ \hline
RF-8 Publicar & CU-005 & \texttt{POST/PUT/DELETE /EventoServlet} y \texttt{/ActividadServlet} & \texttt{EventoServlet, ActividadServlet / ManagerEventos, ActividadManager} & Impl. \\ \hline
RF-9 Panel & CU-011 & \texttt{GET /DashboardServlet} \& \texttt{GET /AuditoriaServlet} & \texttt{DashboardServlet, AuditoriaServlet / ManagerUsuarios, ManagerAuditoria} & Impl. \\ \hline
RF-10 Inscri. & CU-004 & \texttt{POST /InscripcionServlet} \& \texttt{DELETE /InscripcionServlet} & \texttt{InscripcionServlet / ManagerInscripciones} & Impl. \\ \hline
RF-11 Perfil & CU-003 & \texttt{POST /PerfilServlet} \& \texttt{PUT /VisibilidadServlet} & \texttt{PerfilServlet, VisibilidadServlet / AlumniManager} & Impl. \\ \hline
RF-12 Inact. & Proceso Job & \texttt{POST /InactividadServlet} + \texttt{JobsContextListener} (auto 24 h) & \texttt{InactividadServlet, JobsContextListener / ManagerUsuarios} & Impl. \\ \hline
RF-13 Privac. & Transversal & \texttt{GET /PerfilServlet?usuario=\{id\}} & \texttt{PerfilServlet / AlumniManager} & Impl. \\ \hline
RF-14 SMS impl. & CU-012 & \texttt{Intent.ACTION\_SENDTO} (Android) tras \texttt{POST /UsuarioAdminServlet} & \texttt{AdminPanelActivity (cliente Android)} & Impl. \\ \hline
\end{longtable}
\end{center}

% CASOS DE USO
\section*{CU-001: Acceder al Sistema}
\begin{usecasetable}
\textbf{Versión} & 1.1 (29/04/2026) \\
\hline
\textbf{Dependencias} & \dependencies{\item RG-1. \item RF-1. \item RN-1.} \\
\hline
\textbf{Precondición} & El usuario no está autenticado, pero posee credenciales activas en el sistema. \\
\hline
\textbf{Descripción} & Validación de credenciales y generación de sesión segura. \\
\hline
\textbf{Secuencia normal} &
\begin{flowtable}{0.6\textwidth}
1 & Usuario solicita acceder al login e introduce usuario y contraseña. \\
2 & Sistema valida credenciales, comprueba estado activo y rol. \\
3 & Sistema genera token/sesión y redirige a la interfaz correspondiente. \\
\end{flowtable} \\
\hline
\textbf{Postcondición} & El usuario queda autenticado en el sistema con su rol específico identificado. \\
\hline
\textbf{Excepciones} &
\begin{flowtable}{0.6\textwidth}
2 & Credenciales incorrectas o cuenta inactiva: Sistema muestra error 401 y rechaza acceso. \\
\end{flowtable} \\
\hline
\end{usecasetable}

\section*{CU-002: Buscar Alumni}
\begin{usecasetable}
\textbf{Versión} & 1.1 (29/04/2026) \\
\hline
\textbf{Dependencias} & \dependencies{\item RG-2. \item RF-3. \item RN-4.} \\
\hline
\textbf{Precondición} & Usuario autenticado con rol autorizado. \\
\hline
\textbf{Descripción} & Búsqueda simple de alumni respetando privacidad. \\
\hline
\textbf{Secuencia normal} &
\begin{flowtable}{0.6\textwidth}
1 & Usuario introduce un término simple en el buscador. \\
2 & Sistema filtra alumni, eliminando datos privados según RN-4. \\
3 & Sistema devuelve lista de perfiles públicos. \\
\end{flowtable} \\
\hline
\textbf{Postcondición} & Resultados mostrados aplicando filtros de privacidad. \\
\hline
\textbf{Excepciones} &
\begin{flowtable}{0.6\textwidth}
2 & Sin resultados: Mensaje informativo. \\
1 & Consulta vacía masiva: Sistema requiere al menos 3 caracteres. \\
\end{flowtable} \\
\hline
\end{usecasetable}

\section*{CU-003: Modificar Perfil Personal}
\begin{usecasetable}
\textbf{Versión} & 1.1 (29/04/2026) \\
\hline
\textbf{Dependencias} & \dependencies{\item RF-11. \item RN-2.} \\
\hline
\textbf{Precondición} & Usuario autenticado solicita modificar SU propio perfil. \\
\hline
\textbf{Secuencia normal} &
\begin{flowtable}{0.6\textwidth}
1 & Usuario modifica datos y ajustes de privacidad. \\
2 & Sistema valida formato y guarda preferencias de visibilidad. \\
\end{flowtable} \\
\hline
\textbf{Postcondición} & Perfil y preferencias de privacidad actualizadas. \\
\hline
\textbf{Excepciones} &
\begin{flowtable}{0.6\textwidth}
2 & Formato erróneo: Mensaje descriptivo de error. \\
\end{flowtable} \\
\hline
\end{usecasetable}

\section*{CU-004: Inscribirse a Evento/Actividad}
\begin{usecasetable}
\textbf{Versión} & 1.1 (29/04/2026) \\
\hline
\textbf{Dependencias} & \dependencies{\item RF-10. \item RN-6. \item RN-7. \item RN-8. \item RN-9.} \\
\hline
\textbf{Precondición} & Usuario autenticado con rol Alumni o PDI. Evento o actividad aprobado/a, visible, no cancelado/a, dentro del plazo de inscripción y con plazas disponibles. \\
\hline
\textbf{Secuencia normal} &
\begin{flowtable}{0.6\textwidth}
1 & Usuario solicita inscripción. \\
2 & Sistema valida fecha límite, plazas disponibles y no duplicidad. \\
3 & Sistema registra inscripción y actualiza aforo. \\
\end{flowtable} \\
\hline
\textbf{Postcondición} & Usuario registrado. Aforo reducido. \\
\hline
\textbf{Excepciones} &
\begin{flowtable}{0.6\textwidth}
2 & Recurso no aprobado, cancelado, sin plazas, fuera de plazo, usuario ya inscrito o rol no autorizado: el sistema bloquea la acción y muestra el motivo. \\ \hline
\end{flowtable} \\
\hline
\end{usecasetable}

\section*{CU-005: Publicar Evento o Actividad Aprobada}
\begin{usecasetable}
\textbf{Versión} & 1.2 (17/05/2026) \\
\hline
\textbf{Dependencias} & \dependencies{\item RG-3. \item RF-8. \item RN-6. \item CU-010.} \\
\hline
\textbf{Precondición} & Existe una propuesta de evento o actividad evaluada como aprobada. \\
\hline
\textbf{Descripción} & El sistema permite al PTGAS o Administrador formalizar y hacer público en la cartelera general un evento o actividad que previamente ha sido aprobado/a en la fase de revisión. \\
\hline
\textbf{Secuencia normal} &
\begin{flowtable}{0.6\textwidth}
1 & El usuario accede al listado de propuestas aprobadas pendientes de publicación. \\
2 & El usuario selecciona la propuesta aprobada y solicita su publicación definitiva. \\
3 & El sistema valida los datos del recurso definitivo, genera la entidad pública correspondiente en estado activo y la añade a la cartelera general. \\
4 & El sistema notifica automáticamente al Alumni/PDI proponente que su evento o actividad ha sido publicado/a mediante el endpoint \texttt{PUT /PropuestaServlet} con decisión \texttt{PUBLICAR}, creando una entrada en la tabla \texttt{notificacion}. \\
\hline
\end{flowtable} \\
\hline
\textbf{Postcondición} & El evento o actividad pasa a estar visible y disponible para recibir inscripciones (RN-6). \\
\hline
\textbf{Excepciones} &
\begin{flowtable}{0.6\textwidth}
3 & Si la fecha actual es posterior a la fecha del evento: El sistema muestra un error de validación temporal y deniega la publicación. \\
\hline
\end{flowtable} \\
\hline
\end{usecasetable}

\section*{CU-006: Eliminar Cuenta de Usuario}
\begin{usecasetable}
\textbf{Versión} & 1.1 (29/04/2026) \\
\hline
\textbf{Dependencias} & \dependencies{\item RF-6. \item RN-10.  \item RN-11. \item RI-8.} \\
\hline
\textbf{Precondición} & Administrador autenticado. \\
\hline
\textbf{Secuencia normal} &
\begin{flowtable}{0.6\textwidth}
1 & Admin selecciona eliminar cuenta y confirma. \\
2 & Sistema realiza baja lógica, anonimizando datos personales y manteniendo históricos de asistencia huérfanos. Registra en auditoría. \\
\end{flowtable} \\
\hline
\textbf{Postcondición} & Cuenta inactiva/anonimizada. \\
\hline
\textbf{Excepciones} &
\begin{flowtable}{0.6\textwidth}
1 & Operación cancelada por Admin. \\
\end{flowtable} \\
\hline
\end{usecasetable}

\section*{CU-007: Proponer Nuevo Evento o Actividad}
\begin{usecasetable}
\textbf{Versión} & 1.1 (29/04/2026) \\
\hline
\textbf{Dependencias} & \dependencies{\item RG-3. \item RF-4. \item RN-5. \item RI-6.} \\ \\
\hline
\textbf{Precondición} & Usuario Alumni o PDI. \\
\hline
\textbf{Secuencia normal} &
\begin{flowtable}{0.6\textwidth}
1 & Usuario envía formulario de propuesta de evento o actividad. \\
2 & Sistema persiste la propuesta en la base de datos en estado pendiente y notifica a PTGAS. \\
\end{flowtable} \\
\hline
\textbf{Postcondición} & Propuesta registrada (NO visible al público). \\
\hline
\textbf{Excepciones} &
\begin{flowtable}{0.6\textwidth}
1 & Faltan campos clave: Bloqueo de envío. \\
\end{flowtable} \\
\hline
\end{usecasetable}

\section*{CU-008: Buscar Alumni (Filtros)}
\begin{usecasetable}
\textbf{Versión} & 1.1 (29/04/2026) \\
\hline
\textbf{Dependencias} & \dependencies{\item RF-3.   \item RN-4.  \item RF-13.} \\
\hline
\textbf{Precondición} & Usuario autenticado. \\
\hline
\textbf{Secuencia normal} &
\begin{flowtable}{0.6\textwidth}
1 & Usuario introduce 1 o más filtros opcionales. \\
2 & Sistema cruza datos respetando privacidad y devuelve lista. \\
\end{flowtable} \\
\hline
\textbf{Postcondición} & Lista filtrada mostrada. \\
\hline
\textbf{Excepciones} &
\begin{flowtable}{0.6\textwidth}
1 & Si TODOS los filtros están vacíos, se exige al menos uno. \\
\end{flowtable} \\
\hline
\end{usecasetable}

\section*{CU-009: Activar Cuenta de Usuario}
\begin{usecasetable}
\textbf{Versión} & 1.0 (17/05/2026) \\
\hline
\textbf{Dependencias} & \dependencies{\item RG-1. \item RF-2. \item RI-3.} \\
\hline
\textbf{Precondición} & El usuario ha recibido un correo electrónico de la universidad con un enlace seguro y un token temporal. Su cuenta está pendiente de activación. \\
\hline
\textbf{Descripción} & El sistema valida el token temporal de activación, obliga al usuario a definir una contraseña segura que cumpla con las políticas del sistema y activa la cuenta. \\
\hline
\textbf{Secuencia normal} &
\begin{flowtable}{0.6\textwidth}
1 & El usuario accede a la plataforma a través del enlace de activación recibido por correo. \\
2 & El sistema valida que el token no haya expirado y muestra el formulario de asignación de contraseña obligatoria. \\
3 & El usuario introduce su nueva contraseña cumpliendo las políticas de seguridad. \\
4 & El sistema genera el hash salado de la contraseña, actualiza las credenciales, activa el usuario e invalida el token temporal. \\
\hline
\end{flowtable} \\
\hline
\textbf{Postcondición} & La cuenta pasa a estar activa y habilitada para iniciar sesión normal (RN-1). \\
\hline
\textbf{Excepciones} &
\begin{flowtable}{0.6\textwidth}
2 & Token expirado o inválido: El sistema deniega el acceso y solicita al usuario requerir un nuevo enlace de activación. \\
4 & Contraseña débil (no cumple políticas): El sistema muestra los requisitos de seguridad y mantiene al usuario en el formulario. \\
\hline
\end{flowtable} \\
\hline
\end{usecasetable}

\section*{CU-010: Evaluar Propuesta de Evento o Actividad}
\begin{usecasetable}
\textbf{Versión} & 1.1 (29/05/2026) \\
\hline
\textbf{Dependencias} & \dependencies{\item RG-3. \item RF-7. \item RN-5. \item RI-6. \item RI-11.} \\
\hline
\textbf{Precondición} & El usuario es PTGAS o Administrador autenticado. Existen propuestas pendientes en el sistema (RN-5). \\
\hline
\textbf{Descripción} & Permite al personal autorizado de la universidad revisar los detalles de una propuesta de evento o actividad enviada por un Alumni/PDI y emitir un veredicto de aprobación o rechazo justificado. \\
\hline
\textbf{Secuencia normal} &
\begin{flowtable}{0.6\textwidth}
1 & El usuario accede a la bandeja de entrada de propuestas pendientes de evaluación. \\
2 & El sistema muestra la lista de propuestas pendientes. \\
3 & El usuario selecciona una propuesta, revisa los datos (título, descripción, aforo sugerido) y selecciona una acción: aprobar la propuesta o rechazarla. \\
4a & Si selecciona aprobar la propuesta: el sistema cambia el estado de la propuesta para guardarla como aprobada, habilitándola para el CU-005. \\
4b & Si selecciona rechazar la propuesta: el sistema obliga a introducir un motivo de rechazo y la guarda como rechazada. \\
5 & El sistema persiste el cambio de estado en la base de datos. \\
6 & El sistema crea una entrada en la tabla \texttt{notificacion} de tipo \texttt{PROPUESTA\_RESUELTA} dirigida al solicitante mediante \texttt{ManagerNotificaciones} (RI-11). \\
\end{flowtable} \\
\hline
\textbf{Postcondición} & La propuesta cambia de estado en el sistema (RI-6) y sale de la bandeja de pendientes. El solicitante recibe la notificación interna correspondiente. \\
\hline
\textbf{Excepciones} &
\begin{flowtable}{0.6\textwidth}
3 & Propuesta inexistente o no pendiente: el sistema devuelve error 404 y no modifica datos. \\
4b & Rechazo sin motivo: el sistema bloquea la operación y solicita justificación. \\
5 & Error al persistir la decisión: el sistema muestra error descriptivo y no confirma el cambio. \\
\end{flowtable} \\
\hline
\end{usecasetable}

\section*{CU-011: Consultar Panel de Administración y Auditoría}
\begin{usecasetable}
\textbf{Versión} & 1.1 (29/05/2026) \\
\hline
\textbf{Dependencias} & \dependencies{\item RG-7. \item RF-9. \item RI-8. \item RN-11.} \\
\hline
\textbf{Precondición} & El usuario es Administrador autenticado y tiene permisos de acceso al panel de control. \\
\hline
\textbf{Descripción} & Permite al Administrador consultar métricas generales del sistema y revisar el registro de auditoría de acciones administrativas sensibles. \\
\hline
\textbf{Secuencia normal} &
\begin{flowtable}{0.6\textwidth}
1 & El Administrador accede al Panel de Control desde su zona privada. \\
2 & El sistema valida que el usuario autenticado tiene rol Administrador. \\
3 & El sistema solicita las métricas generales mediante \texttt{GET /DashboardServlet}. \\
4 & El sistema muestra usuarios registrados, eventos/actividades publicados, propuestas pendientes, inscripciones y cuentas inactivas. \\
5 & El Administrador accede a la sección de Auditoría. \\
6 & El sistema consulta el registro de acciones sensibles mediante \texttt{GET /AuditoriaServlet}. \\
7 & El sistema muestra el listado de acciones administrativas con fecha, actor, acción realizada y resultado. \\
\end{flowtable} \\
\hline
\textbf{Postcondición} & El Administrador visualiza las métricas y los registros de auditoría sin modificar datos del sistema. \\
\hline
\textbf{Excepciones} &
\begin{flowtable}{0.6\textwidth}
2 & Usuario no autenticado: el sistema devuelve error 401 y deniega el acceso. \\
2 & Usuario sin rol Administrador: el sistema devuelve error 403 y bloquea el acceso al panel. \\
3/6 & Error al cargar métricas o auditoría: el sistema muestra un mensaje descriptivo y registra el fallo. \\
\end{flowtable} \\
\hline
\end{usecasetable}

\section*{CU-012: Crear Cuenta de Usuario con Aviso SMS (Intent Implícito)}
\begin{usecasetable}
\textbf{Versión} & 1.0 (29/05/2026) \\
\hline
\textbf{Dependencias} & \dependencies{\item RG-1. \item RF-6. \item RF-14. \item RI-3.} \\
\hline
\textbf{Precondición} & Administrador autenticado, conectado desde la app Android, y formulario de creación de usuario cumplimentado con un teléfono válido. \\
\hline
\textbf{Descripción} & Permite al administrador dar de alta un nuevo usuario y notificarle por SMS de bienvenida mediante un \textbf{intent implícito} del sistema Android (\texttt{Intent.ACTION\_SENDTO}). El sistema operativo delega la entrega del mensaje en la app de mensajería predeterminada del dispositivo del administrador. \\
\hline
\textbf{Secuencia normal} &
\begin{flowtable}{0.6\textwidth}
1 & El Administrador rellena el formulario de creación de usuario (rol, datos identificativos, teléfono) en \texttt{AdminPanelActivity}. \\
2 & La app envía \texttt{POST /UsuarioAdminServlet} al backend y recibe respuesta de éxito. \\
3 & La app construye un \texttt{Intent} con \texttt{ACTION\_SENDTO} y \texttt{Uri.parse("smsto:" + telefono)}, e incluye el mensaje de bienvenida como extra \texttt{sms\_body}. El mensaje NO contiene la contraseña. \\
4 & La app comprueba mediante \texttt{resolveActivity(getPackageManager())} que existe una app de SMS capaz de atender el intent. \\
5 & La app lanza el intent. El sistema operativo abre la app de mensajería predeterminada con el destinatario y el cuerpo prerrellenados, donde el Administrador confirma el envío. \\
\end{flowtable} \\
\hline
\textbf{Postcondición} & Cuenta creada en el backend y app de SMS abierta en el dispositivo del administrador con el mensaje listo para enviar. \\
\hline
\textbf{Excepciones} &
\begin{flowtable}{0.6\textwidth}
2 & El backend devuelve error de validación: la app muestra el error y no lanza el intent. \\
3 & El teléfono del usuario está vacío: la app muestra Toast informativo y omite el intent. \\
4 & No existe app de SMS instalada (ej. emulador sin Messages): la app muestra Toast informativo y omite el intent. \\
\end{flowtable} \\
\hline
\end{usecasetable}

\section{Diagramas}
\subsection{Diagrama de Clases}
\begin{figure}[H]
    \centering
    \includegraphics[width=1\textwidth]{diagram_final_12.png}
    \caption{Diagrama de Clases del Dominio}
\end{figure}
En esta versión del modelo, los eventos y actividades se representan como recursos gestionables con un flujo común de propuesta, aprobación, publicación e inscripción. Por ello, la entidad principal de gestión puede especializarse o parametrizarse mediante un atributo de tipo de recurso, evitando duplicar clases cuando el comportamiento funcional es equivalente.

\subsection{Diagramas de Secuencia}
\subsubsection{CU-001}
\begin{figure}[H]
    \centering
    \includegraphics[width=1\textwidth]{casu_1.png}
    \label{fig:CU-1}
\end{figure}

\subsubsection{CU-002}
\begin{figure}[H]
    \centering
    \includegraphics[width=1\linewidth]{casu2.png}
    \label{fig:CU-2}
\end{figure}

\subsubsection{CU-003}
\begin{figure}[H]
    \centering
    \includegraphics[width=1\linewidth]{casu3.png}
    \label{fig:CU-3}
\end{figure}

\subsubsection{CU-004}
\begin{figure}[H]
    \centering
    \includegraphics[width=1\linewidth]{casu4.png}
    \label{fig:CU-4}
\end{figure}
\newpage

\subsubsection{CU-005}
\begin{figure}[H]
    \centering
    \includegraphics[width=1\linewidth]{casu5.png}
    \label{fig:CU-5}
\end{figure}
\newpage

\subsubsection{CU-006}
\begin{figure}[H]
    \centering
    \includegraphics[width=1\linewidth]{casu6.png}
    \label{fig:CU-6}
\end{figure}
\newpage

\subsubsection{CU-007}
\begin{figure}[H]
    \centering
    \includegraphics[width=1\linewidth]{casu7.png}
    \label{fig:CU-7}
\end{figure}
\newpage

\subsubsection{CU-008}
\begin{figure}[H]
    \centering
    \includegraphics[width=1\linewidth]{casu8.png}
    \label{fig:CU-8}
\end{figure}

\subsubsection{CU-009}
\begin{figure}[H]
    \centering
    \includegraphics[width=1\linewidth]{cu9_new.png}
    \label{fig:CU-9}
\end{figure}

\subsubsection{CU-010}
\begin{figure}[H]
    \centering
    \includegraphics[width=1\linewidth]{casu10_new.png}
    \label{fig:CU-10}
\end{figure}

\subsubsection{CU-011}
\begin{figure}[H]
    \centering
    \includegraphics[width=1\linewidth]{casu11.png}
    \label{fig:CU-11}
\end{figure}

\subsubsection{CU-012}
\textit{Diagrama de secuencia pendiente de incorporar: PTGAS/Admin $\rightarrow$ \texttt{AdminPanelActivity} $\rightarrow$ \texttt{POST /UsuarioAdminServlet} $\rightarrow$ Backend OK $\rightarrow$ \texttt{Intent.ACTION\_SENDTO} (Android OS) $\rightarrow$ App de mensajería predeterminada (entidad externa). El intent es implícito porque la app no nombra clase destino, sino que delega en el sistema operativo la elección del receptor.}

\subsection{Prototipo}
\begin{figure}[htbp]
    \centering
    \includegraphics[width=0.5\linewidth]{pro1.png}
    \caption{Prototipo inicial (referencia visual previa al desarrollo en Android nativo)}
\end{figure}

\begin{figure}[htbp]
    \centering
    \includegraphics[width=0.5\linewidth]{pro2.png}
    \caption{Prototipo inicial (referencia visual)}
\end{figure}
\clearpage

\begin{figure}[htbp]
    \centering
    \includegraphics[width=0.5\linewidth]{pro3.png}
    \caption{Prototipo inicial (referencia visual)}
\end{figure}

\subsection{Mapa de Navegación}
\begin{figure}[htbp]
    \centering
    \includegraphics[width=0.5\linewidth]{navigation_map.jpeg}
    \caption{Mapa de Navegación}
\end{figure}

\subsubsection*{Descripción de Flujos de Navegación Desarrollados por Rol}
Para solventar la separación de espacios y el ciclo de vida de los eventos, el Mapa de Navegación se estructura bajo un control de acceso basado en roles que cierra el flujo completo del sistema:
\begin{enumerate}
    \item \textbf{Flujo del Proceso de Eventos (Cierre del Ciclo):}
    El ciclo nace en el espacio del \textbf{Alumni o PDI}, quienes navegan a \textit{"Proponer Evento"}, rellenando el formulario que invoca al endpoint \texttt{POST /PropuestaServlet} (\textbf{CU-007}). El evento queda bloqueado al público en estado pendiente. Posteriormente, el \textbf{PTGAS} visualiza esta entrada en su bandeja exclusiva de \textit{"Revisión de Propuestas"}, ejecutando el \textbf{CU-010} mediante el endpoint \texttt{PUT /PropuestaServlet}. Si es aprobado, el PTGAS tiene acceso al botón de acción \textit{"Publicar"}, que transforma la propuesta en un objeto \texttt{Evento} activo en la cartelera general (\textbf{CU-005}), permitiendo que se abra el flujo de inscripción (\textbf{CU-004}) para los Alumni. Las propuestas gestionadas por estos servicios pueden ser de tipo evento o actividad, diferenciándose mediante un atributo interno \texttt{tipoRecurso}.

    \item \textbf{Flujo Exclusivo del Administrador (Admin Loop):}
    Tras autenticarse, el Administrador es redirigido a una pasarela restringida que contiene enlaces directos a dos áreas inaccesibles para el resto de roles: el \textit{"Panel de Control / Dashboard"} (\textbf{CU-011}), que renderiza métricas de inactividad (\texttt{GET /DashboardServlet}), y la consola de \textit{"Auditoría"}, que lee en tiempo real el archivo histórico de acciones administrativas sensibles (\texttt{GET /AuditoriaServlet}), garantizando el control sobre las bajas lógicas (\textbf{CU-006}).

    \item \textbf{Flujo del Personal Docente e Investigador (PDI):}
    Su navegación está estrictamente orientada al networking académico. Desde su página de inicio, su flujo se ramifica únicamente hacia el \textit{"Buscador Avanzado"} (para localizar antiguos alumnos y proponerles colaboraciones o colaboraciones con empresas) y hacia la \textit{"Propuesta de Eventos Co-organizados"}, careciendo de botones de edición sobre perfiles ajenos o publicación directa.

    \item \textbf{Flujo del Administrador con SMS Implícito (Alta de Usuarios):}
    Desde \texttt{AdminPanelActivity}, tras cumplimentar el formulario de creación de usuario y recibir respuesta exitosa del endpoint \texttt{POST /UsuarioAdminServlet}, la app Android dispara un \textbf{intent implícito} (\texttt{Intent.ACTION\_SENDTO} con URI \texttt{smsto:}) que delega en la app de mensajería predeterminada del dispositivo del administrador para enviar un mensaje de bienvenida al nuevo usuario. Este flujo conecta el sistema con una aplicación externa al proyecto a través del bus de intents del sistema operativo (\textbf{CU-012}, \textbf{RF-14}).
\end{enumerate}

\subsubsection*{Matriz de Justificación del Prototipo (Pantalla / Activity $\rightarrow$ Rol $\rightarrow$ Requisito)}
\begin{center}
\small
\renewcommand{\arraystretch}{1.25}
\begin{longtable}{|p{0.24\textwidth}|p{0.25\textwidth}|p{0.43\textwidth}|}
\hline
\textbf{Pantalla / Activity Android} &
\textbf{Roles con Acceso} &
\textbf{Requisitos Funcionales Vinculados} \\
\hline
\endfirsthead
\hline
\textbf{Pantalla / Activity Android} &
\textbf{Roles con Acceso} &
\textbf{Requisitos Funcionales Vinculados} \\
\hline
\endhead
\textbf{LoginActivity} (Vista de Login y Activación de Cuenta) &
Usuario no autenticado, Alumni, PDI, PTGAS y Administrador. &
\textbf{RF-1:} Permite el acceso mediante credenciales y validación de rol.
\newline
\textbf{RF-2:} Permite la activación inicial mediante enlace temporal y cambio obligatorio de contraseña. \\
\hline
\textbf{DirectorioActivity} (Buscador Simple y Avanzado de Alumni) &
Alumni, PDI, PTGAS y Administrador. &
\textbf{RF-3:} Permite realizar búsquedas simples y avanzadas mediante filtros cruzados.
\newline
\textbf{RF-5:} Permite al PTGAS buscar perfiles de Alumni.
\newline
\textbf{RF-13:} Muestra únicamente la información autorizada según rol y preferencias de privacidad. \\
\hline
\textbf{MiPerfilActivity} (Mi Perfil y Configuración de Privacidad) &
Alumni, PDI y PTGAS para modificación de su propio perfil. &
\textbf{RF-5:} Permite al PTGAS modificar su propio perfil.
\newline
\textbf{RF-11:} Permite al Alumni editar su perfil y configurar la visibilidad campo a campo.
\newline
\textbf{RF-13:} Aplica las preferencias de privacidad en la visualización pública del perfil. \\
\hline
\textbf{EventosActivity / ActividadesActivity} (Listado de Eventos y Actividades) &
Alumni, PDI, PTGAS y Administrador. &
\textbf{RF-8:} Muestra eventos aprobados, publicados o cancelados por roles autorizados.
\newline
\textbf{RF-10:} Permite consultar eventos y actividades disponibles para inscripción. \\
\hline
\textbf{DetalleEventoActivity / InscripcionesActivity} (Inscripción y Cancelación de Asistencia) &
Alumni y PDI. &
\textbf{RF-10:} Permite inscribirse o cancelar la inscripción en eventos/actividades aprobadas, siempre que existan plazas y el plazo esté abierto. \\
\hline
\textbf{ProponerActivity} (Formulario de Propuesta de Evento o Actividad) &
Alumni y PDI. &
\textbf{RF-4:} Permite registrar propuestas de eventos y actividades en estado pendiente. \\
\hline
\textbf{PropuestasActivity} (Bandeja de Evaluación de Propuestas) &
PTGAS y Administrador. &
\textbf{RF-7:} Permite revisar, aprobar o rechazar propuestas de eventos o actividades.
\newline
\textbf{RF-8:} Habilita la publicación posterior únicamente de eventos o actividades aprobados. \\
\hline
\textbf{PropuestasActivity (acción Publicar)} (Publicación de Evento Aprobado) &
PTGAS y Administrador. &
\textbf{RF-8:} Permite publicar, modificar o cancelar eventos previamente aprobados.
\newline
\textbf{RF-7:} Mantiene la trazabilidad con la evaluación previa de la propuesta. \\
\hline
\textbf{AdminPanelActivity} (Panel de Administración y Auditoría) &
Administrador. &
\textbf{RF-6:} Permite crear, modificar, suspender y cambiar roles de cuentas.
\newline
\textbf{RF-9:} Muestra métricas de usuarios, eventos, propuestas pendientes, inscripciones y auditoría de acciones.
\newline
\textbf{RF-14:} Lanza el intent implícito de SMS de bienvenida tras crear un usuario. \\
\hline
\textbf{InactividadServlet + JobsContextListener} (Recordatorios e Inactividad) &
Administrador y proceso automático del sistema. &
\textbf{RF-12:} Permite controlar o ejecutar el proceso automático de notificación a cuentas con inactividad superior a un año. El scheduler interno se dispara automáticamente cada 24 h. \\
\hline
\end{longtable}
\end{center}

\section{Diseño del Backend}
En esta sección se detalla la implementación técnica del servidor y del cliente Android, describiendo la arquitectura, los servicios y los patrones de diseño aplicados.

\subsection{Arquitectura del Sistema y Descripción del Backend}
El backend de la aplicación \textit{Desarrollo Alumni} se ha desarrollado siguiendo una arquitectura web clásica por capas. A nivel de infraestructura técnica, el proyecto se especifica de la siguiente manera:
\begin{itemize}
    \item \textbf{Entorno y Servidor:} Desarrollado en \textbf{Java 8} (\texttt{maven.compiler.source/target = 1.8}) bajo el estándar Java Enterprise Edition (Java EE) con la API \texttt{javax.servlet} 3.1. El servidor de aplicaciones objetivo para el despliegue es \textbf{Apache Tomcat 9}, configurado sobre \textbf{XAMPP}.
    \item \textbf{Gestor de Dependencias:} \textbf{Maven}, utilizando un arquetipo \texttt{webapp}. Las dependencias principales incluyen \texttt{javax.servlet-api}, \texttt{org.json}, el conector \textbf{JDBC} para MariaDB/MySQL y \texttt{jBCrypt 0.4} para el hashing de contraseñas, manteniendo un backend ligero.
    \item \textbf{Estructura de Paquetes:}
    \begin{itemize}
        \item \texttt{es.loyola.classes}: Entidades del Modelo de Dominio (\texttt{Usuario}, \texttt{Alumni}, \texttt{Evento}, \texttt{Notificacion}, etc.).
        \item \texttt{es.loyola.enums}: Enumerados de apoyo (roles, facultades, estados).
        \item \texttt{es.loyola.dao}: Capa de acceso a datos e interfaces (\texttt{ManagerUsuarios}, \texttt{AlumniManager}, \texttt{ManagerNotificaciones}, etc.).
        \item \texttt{es.loyola.servlets}: Capa de Controladores (Endpoints).
        \item \texttt{es.loyola.security}: Filtros y utilidades (\texttt{SecurityFilter}, \texttt{PasswordUtil}, \texttt{SessionContext}).
        \item \texttt{es.loyola.jobs}: Tareas programadas internas del backend (\texttt{JobsContextListener}).
    \end{itemize}
    \item \textbf{Persistencia y Base de Datos:} La aplicación incorpora persistencia relacional mediante \textbf{MariaDB/MySQL en XAMPP}, configurada a través de una clase auxiliar \texttt{Database} (singleton) en el paquete \texttt{es.loyola.db}. Las clases \texttt{Manager} del paquete \texttt{es.loyola.dao} actúan como capa DAO de la aplicación, encapsulando las operaciones SQL contra la base de datos. De este modo, los \texttt{Servlets} mantienen su responsabilidad como controladores HTTP, delegando la lógica de acceso a datos en los Managers/DAOs. Este enfoque sustituye la persistencia en memoria por una solución persistente, evitando la pérdida de datos tras reinicios y permitiendo una evolución posterior hacia una arquitectura más robusta con JDBC estructurado o JPA/Hibernate, sin alterar la firma pública de los \texttt{Servlets}.
\end{itemize}

\subsubsection*{Cliente Android nativo}
El cliente principal del sistema es una \textbf{app Android nativa} (no una PWA) ubicada en el módulo \texttt{Alumni\_android}, paquete raíz \texttt{es.uloyola.pdm.Alumni\_android}. Características técnicas relevantes:
\begin{itemize}
    \item \textbf{Lenguaje y SDK:} Java 11, \texttt{compileSdk 36}, \texttt{minSdk 24}, AGP 9.1.1, Gradle KTS.
    \item \textbf{Capa de red:} \texttt{Retrofit 2.9.0} sobre \texttt{OkHttp 4.9.3}, con conversor \texttt{Gson}. Cliente centralizado en \texttt{conexionServer/RetrofitClient.java} apuntando a \texttt{https://10.0.2.2:8443/alumni/} (IP especial del emulador hacia el host).
    \item \textbf{Sesión:} \texttt{AlumniCookieJar} persiste la cookie \texttt{JSESSIONID} entre llamadas. Singleton \texttt{SessionManager} mantiene el usuario y rol activos en memoria del cliente.
    \item \textbf{Seguridad de transporte:} \texttt{network\_security\_config.xml} confía en el certificado autofirmado del backend (\texttt{res/raw/alumni\_cert.crt}) para permitir HTTPS / TLS 1.3 en entorno local.
    \item \textbf{Patrón de pantalla:} cada vista es una \texttt{Activity}; navegación interna mediante \textbf{intents explícitos} (\texttt{new Intent(this, OtraActivity.class)}) y un \textbf{intent implícito} (\texttt{ACTION\_SENDTO} con \texttt{smsto:}) en \texttt{AdminPanelActivity} para el SMS de bienvenida (RF-14).
\end{itemize}

\subsubsection*{Copias de seguridad (RNF-4)}
La estrategia de copias se materializa como cuatro scripts \texttt{.bat} en la carpeta \texttt{alumni/backup/} del proyecto:
\begin{itemize}
    \item \texttt{backup\_completo.bat} -- volcado total semanal con \texttt{mysqldump}, cifrado AES-256 con \texttt{openssl} y duplicado en dos rutas (\texttt{C:\textbackslash backups\_alumni\textbackslash completas} y \texttt{D:\textbackslash backups\_alumni\_redundante\textbackslash completas}), purga > 60 días.
    \item \texttt{backup\_incremental.bat} -- equivalente para copias diarias.
    \item \texttt{programar\_tareas.bat} -- registra ambos en el Programador de Tareas de Windows (\texttt{schtasks}): incremental diaria a las 03:00 y completa los domingos a las 03:30 (dentro de la ventana valle 02:00--05:00 del RNF-3).
    \item \texttt{restaurar\_copia.bat} -- procedimiento de recuperación desde una copia cifrada.
\end{itemize}

\subsubsection*{Tareas programadas internas (RF-12)}
Además de la herramienta externa de copias, el backend implementa un \textbf{scheduler interno} en el paquete \texttt{es.loyola.jobs}:
\begin{itemize}
    \item \texttt{JobsContextListener}: \texttt{@WebListener} que implementa \texttt{ServletContextListener} y arranca con el contexto de Tomcat. Usa un \texttt{ScheduledExecutorService} con hilo daemon que se arma al despliegue y se detiene en \texttt{contextDestroyed}.
    \item \textbf{Job principal:} recordatorio de inactividad. Calcula los segundos hasta la próxima 03:00 local y programa una ejecución periódica cada 24 horas. Para cada cuenta con inactividad superior a 365 días, crea una entrada en \texttt{notificacion} (tipo \texttt{INACTIVIDAD}) y registra el ciclo completo en auditoría.
\end{itemize}
El sistema combina por tanto dos esquemas de programación: un \textbf{scheduler externo} basado en \texttt{schtasks} (para copias de seguridad) y un \textbf{scheduler interno} basado en \texttt{ScheduledExecutorService} (para notificaciones e inactividad).

El siguiente diagrama ilustra la arquitectura de capas, incluyendo el cliente Android, el backend en capas y la integración futura:
\begin{figure}[H]
    \centering
    \includegraphics[width=0.7\textwidth]{diagram-export-5-17-2026-4_21_36-PM.png}
    \caption{Diagrama de Arquitectura de Capas y Futura Integración}
\end{figure}

\subsection{Descripción del Diagrama de Clases y Multiplicidades}
Las clases \texttt{Usuario}, \texttt{Alumni}, \texttt{Pdi}, \texttt{Ptgas}, \texttt{Administrador}, \texttt{Evento}, \texttt{Actividad}, \texttt{Inscripcion}, \texttt{Credenciales}, \texttt{Trabajo}, \texttt{Hobbie}, \texttt{Organizador}, \texttt{PropuestaEvento}, \texttt{Notificacion} y \texttt{RegistroAuditoria} están implementadas en el paquete \texttt{es.loyola.classes}. Solo la entidad \texttt{IntegracionSalesforce} se documenta como entidad conceptual, prevista para una integración futura.
\begin{itemize}
    \item \textbf{Jerarquía de Usuario:} \texttt{Usuario} (Clase base abstracta) de la que heredan \texttt{Alumni}, \texttt{PDI}, \texttt{PTGAS} y \texttt{Administrador}. Todos tienen una relación [1..1] con la clase \texttt{Credenciales} (que custodia el hash de contraseña y el token de activación).
    \item \textbf{Propuestas y Aprobaciones:} La clase \texttt{PropuestaEvento} separa la sugerencia de la publicación. Un \texttt{Usuario} (Alumni/PDI) puede tener [0..N] propuestas. Contiene un atributo de tipo \texttt{EstadoPropuesta} (ENUM: PENDIENTE, APROBADA, RECHAZADA, PUBLICADA) y un atributo \texttt{tipoRecurso} que distingue entre EVENTO y ACTIVIDAD.
    \item \textbf{Relación Eventos e Inscripciones:} Un \texttt{Evento} o \texttt{Actividad} puede tener [0..N] \texttt{Usuarios} inscritos a través de la clase intermedia \texttt{Inscripcion} [1..1] Usuario, [1..1] Evento, conteniendo la fecha de registro y la asistencia.
    \item \textbf{Privacidad y Seguridad:} Cada \texttt{Alumni} dispone de una matriz de preferencias de visibilidad campo a campo (\texttt{mostrarContacto}, \texttt{mostrarEmail}, etc.), gestionada por \texttt{AlumniManager} y expuesta por \texttt{VisibilidadServlet}. La entidad \texttt{RegistroAuditoria} acumula [0..N] entradas por cada acción sensible (eliminar cuenta, evaluar propuesta, etc.).
    \item \textbf{Notificaciones (RI-11):} La entidad \texttt{Notificacion} es real e implementada (tabla \texttt{notificacion} en el esquema). Se asocia [0..N] con \texttt{Usuario} (destinatario) y se crea desde \texttt{PropuestaServlet} (notificación al solicitante al aprobar/rechazar/publicar) y desde \texttt{JobsContextListener} (notificación de inactividad).
    \item \textbf{Hobbies y Trabajo:} \texttt{Hobbie} se asocia N:M con \texttt{Alumni} mediante la tabla pivote \texttt{AlumniHobbie}, que incluye el atributo \texttt{experiencia} de tipo \texttt{ExperienciaHobbies} (PRINCIPIANTE, INTERMEDIO, EXPERTO). \texttt{Alumni} mantiene un historial 1..0..* con \texttt{Trabajo}.
    \item \textbf{Organizadores:} \texttt{Organizador} (nombre, código) se asocia con \texttt{Evento} para identificar al responsable institucional de un recurso publicado.
    \item \textbf{Integración futura:} La interfaz \texttt{IntegracionSalesforce} se mantiene como entidad conceptual prevista para la sincronización de perfiles Alumni [1..1] con el CRM externo (no implementada).
\end{itemize}

\subsection{Catálogo de Servicios API (Servlets)}
A continuación se describen todos los servicios que conforman la API del sistema con sus rutas reales tal como están desplegadas en Tomcat bajo el contexto \texttt{/alumni}, incluyendo aquellos proyectados para integraciones futuras.
\begin{longtable}{|p{0.32\textwidth}|p{0.28\textwidth}|p{0.35\textwidth}|}
\hline
\textbf{Servicio (Método y Ruta)} & \textbf{Precondiciones / Rol} & \textbf{Postcondiciones} \\ \hline
\endfirsthead
\hline
\textbf{Servicio (Método y Ruta)} & \textbf{Precondiciones / Rol} & \textbf{Postcondiciones} \\ \hline
\endhead
\multicolumn{3}{|c|}{\textbf{Autenticación}} \\ \hline
\texttt{POST /LoginServlet} & Todos. Cuenta activa, usuario y password válidos. & Inicia sesión, devuelve rol. Error 401 si falla. \\ \hline
\texttt{POST /LogoutServlet} & Usuario autenticado. & Cierra sesión e invalida la cookie JSESSIONID. \\ \hline
\texttt{POST /ActivarCuentaServlet} & Usuario con token temporal válido. Cuenta pendiente de activación. & Activa la cuenta, guarda la nueva contraseña con hash bcrypt e invalida el token temporal. Devuelve error si el token es inválido o ha expirado. \\ \hline
\multicolumn{3}{|c|}{\textbf{Perfiles y Privacidad}} \\ \hline
\texttt{GET /BuscarAlumniServlet} & Autenticados autorizados. & Búsqueda con filtros, respetando privacidad. \\ \hline
\texttt{GET /PerfilServlet?usuario=\{id\}} & Usuarios autorizados. & Devuelve solo campos visibles del perfil. \\ \hline
\texttt{GET /PerfilServlet} & Usuario autenticado. & Devuelve y consulta su propio perfil. \\ \hline
\texttt{POST /PerfilServlet} & Usuario autenticado (o ADMIN). & Edita el perfil con datos válidos. \\ \hline
\texttt{GET /VisibilidadServlet} & Alumni autenticado. & Devuelve las preferencias de privacidad actuales campo a campo. \\ \hline
\texttt{PUT /VisibilidadServlet} & Alumni autenticado. & Actualiza preferencias de visibilidad campo a campo. \\ \hline
\multicolumn{3}{|c|}{\textbf{Eventos y Propuestas}} \\ \hline
\texttt{GET /EventoServlet} & Usuarios autenticados. & Lista o detalle de eventos aprobados. \\ \hline
\texttt{POST/PUT/DELETE /EventoServlet} & PTGAS o Administrador autenticado; recurso/propuesta existente y previamente aprobada cuando la acción sea publicación. & El recurso queda publicado y visible solo después de validar que procede de una propuesta aprobada. \\ \hline
\texttt{POST /PropuestaServlet} & Alumni, PDI autenticados. & Envía propuesta para evaluación. \\ \hline
\texttt{GET /PropuestaServlet} & PTGAS, Admin (o el propio solicitante para las suyas). & Lista propuestas para revisión. \\ \hline
\texttt{PUT /PropuestaServlet} & PTGAS, Admin. & Aprueba, rechaza o publica propuestas; crea automáticamente una \texttt{Notificacion} dirigida al solicitante. \\ \hline
\multicolumn{3}{|c|}{\textbf{Actividades}} \\ \hline
\texttt{GET/POST/PUT/DELETE /ActividadServlet} & PTGAS o Administrador autenticado. Para publicar, debe existir una propuesta de actividad previamente aprobada mediante CU-010. & Publica únicamente actividades aprobadas procedentes del flujo de propuesta/evaluación; permite consultar, modificar o cancelar actividades ya publicadas. \\ \hline
\multicolumn{3}{|c|}{\textbf{Inscripciones}} \\ \hline
\texttt{POST /InscripcionServlet} & Alumni/PDI autenticado. Recurso evento/actividad existente, aprobado, visible, no cancelado, dentro de plazo, con plazas disponibles y sin inscripción previa del usuario. & Registra la inscripción, actualiza el aforo disponible y devuelve confirmación. \\ \hline
\texttt{DELETE /InscripcionServlet} & Usuario previamente inscrito. & Cancela inscripción si la política lo permite. \\ \hline
\multicolumn{3}{|c|}{\textbf{Administración y Auditoría}} \\ \hline
\texttt{POST/GET/PUT /UsuarioAdminServlet} & Administrador autenticado. Usuario objetivo existente cuando proceda. & Crea, consulta, modifica, suspende o cambia roles de usuarios. En operaciones de baja lógica, anonimiza datos personales, conserva históricos no identificables y registra la acción en auditoría. \\ \hline
\texttt{GET /DashboardServlet} & Admin autenticado. & Devuelve métricas (uso, eventos, inactividad). \\ \hline
\texttt{GET /AuditoriaServlet} & Admin autenticado. & Consulta acciones sensibles realizadas. \\ \hline
\multicolumn{3}{|c|}{\textbf{Notificaciones internas (RI-11)}} \\ \hline
\texttt{GET /NotificacionServlet} & Usuario autenticado. & Devuelve la bandeja de notificaciones del usuario. Con \texttt{?unread=1} devuelve solo el conteo de no leídas. \\ \hline
\texttt{PUT /NotificacionServlet} & Usuario autenticado, propietario de la notificación. & Marca una notificación como leída. \\ \hline
\multicolumn{3}{|c|}{\textbf{Recordatorios e Integración}} \\ \hline
\texttt{POST /InactividadServlet} & Administrador. Además se dispara automáticamente cada 24 h vía \texttt{JobsContextListener} (RF-12). & Detecta inactividad y genera recordatorios + notificaciones internas. \\ \hline
\texttt{POST /api/salesforce/sync} & Admin o Servicio interno. \textit{(Endpoint conceptual, no implementado.)} & Sincroniza datos con Salesforce según contrato. \\ \hline
\end{longtable}

\subsection{Patrones de Diseño Utilizados}
\begin{enumerate}
    \item \textbf{Patrón DAO / Repository:} Las clases \texttt{Manager} del paquete \texttt{es.loyola.dao} actúan como capa DAO de la aplicación, encapsulando las operaciones SQL contra la base de datos relacional MariaDB/MySQL. Los \texttt{Servlets} delegan en estas clases el acceso a datos, manteniendo separada la lógica de control HTTP de la persistencia.
    \item \textbf{Gestión centralizada de conexión (Singleton):} La conexión con la base de datos se centraliza mediante una clase auxiliar \texttt{Database} en el paquete \texttt{es.loyola.db}. Esto evita duplicar la lógica de conexión en cada Servlet o Manager.
    \item \textbf{Patrón MVC (Modelo-Vista-Controlador):} Capa \texttt{es.loyola.classes} (Modelo), paquete \texttt{es.loyola.servlets} (Controlador) y la app Android nativa (Vista).
    \item \textbf{Patrón Helper / Utility:} La clase estática \texttt{ServletJsonUtil} centraliza lectura de buffers JSON, formateo de fechas y creación estructurada de respuestas (errores, HTTP codes), evitando duplicidad masiva en Servlets.
    \item \textbf{Patrón Listener / Observer (ciclo de vida del contexto):} \texttt{JobsContextListener} implementa \texttt{ServletContextListener} y se engancha al ciclo de vida de Tomcat para arrancar y parar tareas periódicas con \texttt{ScheduledExecutorService}. Cubre los recordatorios de inactividad (RF-12).
    \item \textbf{Scheduled task externo (Programador del SO):} \texttt{schtasks} de Windows + \texttt{mysqldump} + \texttt{openssl} (AES-256) para las copias de seguridad cifradas y redundantes (RNF-4), independientes del backend Java.
    \item \textbf{Intent implícito (Android):} \texttt{Intent.ACTION\_SENDTO} con URI \texttt{smsto:} en \texttt{AdminPanelActivity} delega en el sistema operativo del dispositivo la entrega del SMS de bienvenida al nuevo usuario, sin nombrar clase destino y sin requerir el permiso peligroso \texttt{SEND\_SMS}. Cubre el RF-14 / CU-012.
    \item \textbf{Filtro de seguridad transversal (Front Controller / Filter Chain):} \texttt{SecurityFilter} intercepta todas las peticiones HTTP, valida sesión y rol, y aplica codificación UTF-8 antes de delegar en el Servlet correspondiente, garantizando el RNF-10.
\end{enumerate}

\section{Glosario y Bibliografía}
\subsection{Glosario}
\begin{itemize}
    \item \textbf{RGPD / Privacidad por Defecto:} Reglamento General de Protección de Datos. Obliga a que los datos solo sean visibles bajo consentimiento explícito.
    \item \textbf{Anonimización:} Proceso de borrado lógico donde se eliminan campos identificativos (nombre, correo) para mantener el histórico de estadísticas (ej. asistencia a eventos) sin identificar al sujeto.
    \item \textbf{PWA (Progressive Web App):} Aplicación web instalable nativamente en móviles, capaz de usar sensores locales. En este proyecto queda como entorno alternativo / futuro.
    \item \textbf{App Android nativa:} Aplicación construida con el SDK de Android, instalable como APK, capaz de usar intents, permisos y sensores del sistema. Cliente principal del proyecto.
    \item \textbf{Intent explícito (Android):} \texttt{Intent} que nombra la clase destino (\texttt{new Intent(this, OtraActivity.class)}). Se usa para la navegación interna entre activities del propio proyecto.
    \item \textbf{Intent implícito (Android):} \texttt{Intent} sin clase destino explícita. Indica una acción (\texttt{ACTION\_SENDTO}, \texttt{ACTION\_VIEW}, etc.) y una URI; el sistema operativo decide qué app del dispositivo lo atiende.
    \item \textbf{\texttt{ACTION\_SENDTO} + URI \texttt{smsto:}}: Combinación que abre la app de mensajería predeterminada con el destinatario y el cuerpo prerrellenados. \textbf{No requiere el permiso \texttt{SEND\_SMS}} porque la app no envía el mensaje, sino que delega la entrega en el sistema operativo y la app de mensajería.
    \item \textbf{Notificación interna:} Entidad \texttt{Notificacion} persistida en la tabla \texttt{notificacion}, gestionada por \texttt{ManagerNotificaciones} y consumida por \texttt{NotificacionServlet}. Distinta de las notificaciones push del SO.
    \item \textbf{Scheduler / \texttt{ServletContextListener}:} Componente Java EE que se engancha al ciclo de vida del contexto del servlet en Tomcat (\texttt{contextInitialized} / \texttt{contextDestroyed}). En este proyecto, \texttt{JobsContextListener} usa este mecanismo junto a un \texttt{ScheduledExecutorService} para disparar tareas periódicas (RF-12).
    \item \textbf{\texttt{ScheduledExecutorService}:} API de \texttt{java.util.concurrent} para programar tareas con retraso o frecuencia fija. Usado por el scheduler interno del backend.
    \item \textbf{AES-256 / \texttt{openssl}:} Algoritmo de cifrado simétrico y herramienta CLI utilizada para cifrar las copias de seguridad en RNF-4 antes de almacenarlas en las ubicaciones redundantes.
    \item \textbf{\texttt{schtasks}:} Utilidad del Programador de Tareas de Windows. En este proyecto se usa para registrar las copias incrementales y completas como tareas automáticas del SO (RNF-4).
    \item \textbf{\texttt{JSESSIONID}:} Cookie generada por Tomcat para mantener la sesión HTTP. Marcada como \texttt{HttpOnly} y \texttt{Secure}; persistida en el cliente Android por \texttt{AlumniCookieJar}.
\end{itemize}

\subsection{Bibliografía}
\begin{itemize}
    \item \textbf{Diseño y prototipado}: Figma. \url{https://www.figma.com/}
    \item \textbf{UML}: PlantUML. \url{https://plantuml.online/}
    \item \textbf{Referencias}: Alumni Global IE. \url{https://www.ie.edu/alumni/} | ESADE Alumni. \url{https://www.esadealumni.net/es}
    \item \textbf{Android}: Guía oficial de Intents. \url{https://developer.android.com/guide/components/intents-filters}
\end{itemize}

\end{document}
